\documentclass[dvipsnames]{article}

\usepackage[utf8]{inputenc}
\usepackage{fontspec}
\usepackage{polyglossia}
\usepackage{unicode-math}

\setdefaultlanguage{russian}
\setotherlanguage{english}

\setmainfont{CMU Serif}
\setsansfont{CMU Sans Serif}
\setmonofont{CMU Typewriter Text}
\newfontfamily\cyrillicfonttt{Menlo}[Scale=MatchLowercase]
% \setmathfont{Latin Modern Math}

\usepackage[letterpaper,top=2cm,bottom=2cm,left=3cm,right=3cm,marginparwidth=1.75cm]{geometry}

\usepackage{amsmath}
\usepackage{fancyhdr}
\usepackage{bussproofs}
\usepackage[colorlinks=true, allcolors=blue]{hyperref}
\usepackage[backend=biber]{biblatex}
\usepackage{xcolor}

\usepackage{minted2}
\usemintedstyle{xcode}
\setminted{escapeinside=||}
\newmintinline[code]{c}{fontsize=\scriptsize}
\let\codeOrig\code
\renewcommand{\code}[2][]{%
  \ifmmode
    \if\relax\detokenize{#1}\relax
      \mathchoice{\mbox{\normalfont\scriptsize\ttfamily #2}}{\mbox{\normalfont\scriptsize\ttfamily #2}}{\mbox{\normalfont\scriptsize\ttfamily #2}}{\mbox{\normalfont\scriptsize\ttfamily #2}}%
    \else
      \mathchoice{\mbox{\normalfont\scriptsize\ttfamily #2}}{\mbox{\normalfont\scriptsize\ttfamily #2}}{\mbox{\normalfont\scriptsize\ttfamily #2}}{\mbox{\normalfont\scriptsize\ttfamily #2}}%
    \fi
  \else
    \codeOrig[#1]{#2}%
  \fi}

% Command for primitive recursion syntax: \primrec{r}{t_zero}{k}{t_nat}{t_step}
\newcommand{\primrec}[5]{%
  \code{with}\,\colorbox{green!20}{\code{#1}}\,\code{:=}\,\colorbox{yellow!30}{\code{#2}}\,\code{for}\,\colorbox{green!20}{\code{#3}}\,\code{=}\,\colorbox{yellow!30}{\code{#4}}\,\code{downto 0 do}\,\colorbox{green!20}{\code{#1}}\,\code{:=}\,\colorbox{yellow!30}{\code{#5}}%
}

\usepackage{mdframed}

\addbibresource{../references.bib}

\begin{document}

\pagestyle{fancy}
\fancyhf{}
\chead{Задание выполнил \textcolor{blue}{Бенжамин Пирс (замените на своё имя)}}
\cfoot{Языки программирования. Семантика и системы типов (весна 2026, Университет ИТМО)}

\section*{Тема 2. Теоретическое задание \textcolor{blue}{(выполнил Бенжамин Пирс)}}%

В следующих заданиях рассматривается расширение простого типизированного $\lambda$-исчисления оператором примитивной рекурсии с синтаксисом, аналогичным циклу \code{for}:
\[
\primrec{r}{t|$_{\mathsf{zero}}$|}{k}{t|$_{\mathsf{nat}}$|}{t|$_{\mathsf{step}}$|}
\]
Семантика: при \code{t|$_{\mathsf{nat}}$|} = \code{0} результат равен \code{t|$_{\mathsf{zero}}$|};
при \code{t|$_{\mathsf{nat}}$|} = \code{succ nv} (где \code{nv} — значение типа \code{Nat}) — в цикле выполняется \code{nv} итераций,
на каждой итерации вычисляется \code{t|$_{\mathsf{step}}$|} с подстановкой текущего значения счётчика \code{k} и текущего значения результата \code{r}, результат этого вычисления становится новым значением \code{r} для следующей итерации; итоговый результат равен значению \code{r} после завершения всех итераций.

\begin{prooftree}
  \AxiomC{$\code{t$_{\mathsf{zero}}$} \longrightarrow \code{t$_{\mathsf{zero}}'$}$}
  \RightLabel{E-NatFor$_1$}
  \UnaryInfC{$\primrec{r}{t|$_{\mathsf{zero}}$|}{k}{t|$_{\mathsf{nat}}$|}{t|$_{\mathsf{step}}$|} \longrightarrow \primrec{r}{t|$_{\mathsf{zero}}'$|}{k}{t|$_{\mathsf{nat}}$|}{t|$_{\mathsf{step}}$|}$}
\end{prooftree}
\begin{prooftree}
  \AxiomC{$\code{t$_{\mathsf{nat}}$} \longrightarrow \code{t$_{\mathsf{nat}}'$}$}
  \RightLabel{E-NatFor$_2$}
  \UnaryInfC{$\primrec{r}{v|$_{\mathsf{zero}}$|}{k}{t|$_{\mathsf{nat}}$|}{t|$_{\mathsf{step}}$|} \longrightarrow \primrec{r}{v|$_{\mathsf{zero}}$|}{k}{t|$_{\mathsf{nat}}'$|}{t|$_{\mathsf{step}}$|}$}
\end{prooftree}
\begin{prooftree}
  \AxiomC{}
  \RightLabel{E-NatForZero}
  \UnaryInfC{$\primrec{r}{v|$_{\mathsf{zero}}$|}{k}{0}{t|$_{\mathsf{step}}$|} \longrightarrow \code{v$_{\mathsf{zero}}$}$}
\end{prooftree}
\begin{prooftree}
  \AxiomC{}
  \RightLabel{E-NatForSucc}
  \UnaryInfC{$\primrec{r}{v|$_{\mathsf{zero}}$|}{k}{succ nv}{t|$_{\mathsf{step}}$|}$}
  \noLine \UnaryInfC{$\longrightarrow \primrec{r}{[k |$\mapsto$| nv, r |$\mapsto$| v|$_{\mathsf{zero}}$|]t|$_{\mathsf{step}}$|}{k}{n}{t|$_{\mathsf{step}}$|}$}
\end{prooftree}

Правило типизации:
\begin{prooftree}
  \AxiomC{$\Gamma \vdash \code{t$_{\mathsf{nat}}$} : \code{Nat}$}
  \AxiomC{$\Gamma \vdash \code{t$_{\mathsf{zero}}$} : \code{T}$}
  \AxiomC{$\Gamma,\, \code{k : Nat},\, \code{r : T} \vdash \code{t$_{\mathsf{step}}$} : \code{T}$}
  \RightLabel{T-NatFor}
  \TrinaryInfC{$\Gamma \vdash \primrec{r}{t|$_{\mathsf{zero}}$|}{k}{t|$_{\mathsf{nat}}$|}{t|$_{\mathsf{step}}$|} : \code{T}$}
\end{prooftree}

\begin{enumerate}

  \item Проставьте явные аннотации типа для аргументов $\lambda$-абстракций и постройте дерево вывода типа для замкнутого терма:
  \[
  \code{(λx. \primrec{r}{λf. f}{k}{pred x}{λf. λx. r (r x)}) (succ 0)}
  \]

  \color{blue}
  \textbf{Решение.}
  \emph{место для решения}
  \color{black}

  \clearpage
  \item
  Расширьте \emph{лемму об инверсии}~\cite[Лемма~9.3.1]{TaPL-ru} и \emph{лемму о канонических формах}~\cite[Лемма~9.3.4]{TaPL-ru} с учётом оператора \code{for}:
  \begin{enumerate}
  \item Добавьте пункт инверсии для термов вида $\primrec{r}{t|$_{\mathsf{zero}}$|}{k}{t|$_{\mathsf{nat}}$|}{t|$_{\mathsf{step}}$|}$: какие типы и выводы обязаны иметь подтермы?
  \item Уточните канонические формы для типа \code{Nat}: остаётся ли утверждение «значение типа \code{Nat} есть либо \code{0}, либо \code{succ nv}»
  верным в расширенном языке, или значения могут иметь и другую форму? Обоснуйте.
  \end{enumerate}

  \color{blue}
  \textbf{Решение.}
  \emph{место для решения}
  \color{black}

  \clearpage
  \item
  Расширьте формулировки и доказательства лемм \emph{продвижения} и \emph{сохранения}~\cite[\S 9.3]{TaPL-ru} на исчисление с оператором \code{for} (примитивная рекурсия).
  Укажите, какие новые случаи появляются в доказательствах и как они разбираются.
  Объясните, почему замкнутый хорошо типизированный терм по-прежнему не может быть тупиковым.

  \color{blue}
  \textbf{Решение.}
  \emph{место для решения}
  \color{black}

  \clearpage
  \item
  Расширьте правила \emph{двунаправленной типизации}~\cite{DunfieldKrishnaswami2021} на оператор \code{for}: введите правила для суждений $\Gamma \vdash t \Leftarrow T$ и $\Gamma \vdash t \Rightarrow T$ (проверка и синтез типа).
  Объясните, почему для \code{for} выбран именно такой режим (проверка, синтез, или оба режима).

  \color{blue}
  \textbf{Решение.}
  \emph{место для решения}
  \color{black}

\end{enumerate}

\section*{Список литературы}
\printbibliography[heading=none]

\end{document}
